% ============================================================
% OP-Q11 CONVERSE CALCULATION — momentum-space OS test
% Version: v2 (2026-06-10). STATUS: DRAFT FOR GPT VERIFICATION.
% Supersedes v1. All review-round corrections implemented:
% no [ordering factor] placeholder; Q_Gr and Q_com defined
% separately; M_Gr derived step by step; sign of m*gamma_0 fixed
% by derivation; symmetrisation shortcut REPLACED by exact
% reordering-sign mechanism; second horn demoted to diagnostic.
% ============================================================

\section*{Momentum-space OS test for the free minimal Euclidean
Dirac kernel (v2)}

% ------------------------------------------------------------
\subsection*{0. Conventions (trial set, complete)}
Euclidean $x=(w,\vec x)$, reflection $\Theta:(w,\vec x)\to(-w,\vec x)$,
consistent with $t=w$.
Hermitian gammas $\{\gamma_A,\gamma_B\}=2\delta_{AB}$,
$\gamma_A^\dagger=\gamma_A$.
Mass gap $m>0$ throughout (assumption of Q11.1a).
Minimal kernel: $D_E=\gamma_A\partial_A+m$,
$\tilde S_E(p)=\dfrac{m-i\gamma_Ap_A}{p^2+m^2}$,
from the Grassmann Gaussian (Berezin) pairing
$\langle\psi_\alpha(x)\bar\psi_\beta(y)\rangle=S_{E,\alpha\beta}(x-y)$.
Reordering rule (THE only point where statistics enters):
\[
\langle\bar\psi_\beta(y)\psi_\alpha(x)\rangle
=\sigma\,S_{E,\alpha\beta}(x-y),
\qquad
\sigma=\begin{cases}-1 & \text{Grassmann},\\
+1 & \text{commuting}.\end{cases}
\]
Antilinear reflection on fields (trial convention; the spinorial
reflection matrix $\gamma_0$ must be checked against the Euclidean
adjoint and charge-conjugation convention adopted in the main text):
\[
\Theta\psi(x)=\bar\psi(\Theta x)\,\gamma_0,
\qquad
\Theta\bar\psi(x)=\gamma_0\,\psi(\Theta x),
\qquad
\Theta(AB)=\Theta(B)\,\Theta(A).
\]
OS sesquilinear form (order fixed; see Remark~0.1):
\[
(F,G)\;\equiv\;\bigl\langle F\cdot\Theta G\bigr\rangle .
\]
Two separate quadratic forms are studied:
$Q_{\rm Gr}[\cdot]=(F,F)$ with Grassmann fields ($\sigma=-1$) and
$Q_{\rm com}[\cdot]=(F,F)$ with commuting spinor fields
($\sigma=+1$), \emph{with the same kernel $S_E$ and the same
reflection data}.

\emph{Remark 0.1 (order convention).}
At the quadratic level the alternative assignment
$(F,G)=\langle\Theta F\cdot G\rangle$ differs from the one above by
the global reordering sign $\sigma$; for Grassmann fields exactly
one of the two assignments is the OS-positive convention.
The v1 draft left this as an unspecified ``ordering factor'' and
thereby produced the $m\gamma_0$ sign mismatch flagged in review;
v2 removes the ambiguity by fixing the order and deriving every
sign.
Notation: $\omega=\sqrt{\vec p^{\,2}+m^2}$,
$h_\pm(\vec p)\equiv m\gamma_0\pm i\gamma_0\vec\gamma\cdot\vec p$.
Both are hermitian
($(i\gamma_0\vec\gamma)^\dagger=-i\vec\gamma\gamma_0
=+i\gamma_0\vec\gamma$) and satisfy
$h_\pm^2=m^2+\vec p^{\,2}=\omega^2$
(cross terms cancel by $\{\gamma_0,\vec\gamma\}=0$),
so $\tfrac{1}{2\omega}(\omega\pm h_\pm)$ are orthogonal rank-2
projectors.

% ------------------------------------------------------------
\subsection*{1. Contour lemma}
For $W>0$ and either sign $s=\pm$:
\[
\int\frac{dp_0}{2\pi}\,
\frac{e^{\,is\,p_0W}\,(m-i\gamma_0p_0-i\vec\gamma\cdot\vec p)}
{p_0^2+\omega^2}
=\frac{e^{-\omega W}}{2\omega}\,
\bigl(m+s\,\omega\gamma_0-i\vec\gamma\cdot\vec p\bigr),
\]
by closing in the upper ($s=+$, pole $p_0=i\omega$) or lower
($s=-$, pole $p_0=-i\omega$) half-plane.
[Check $s=-$: $-i\gamma_0(-i\omega)=-\omega\gamma_0$; prefactor
$(-2\pi i)/(2\pi\cdot(-2i\omega))=1/(2\omega)$.]

% ------------------------------------------------------------
\subsection*{2. Grassmann case: both quadratic families are
reflection positive}
\emph{Family I.}\;
$F=\Psi[f]\equiv\int d^4x\,f^\dagger(x)\,\psi(x)$, $f$ supported in
$w>0$.
Then $\Theta F=\int\bar\psi(\Theta y)\,\gamma_0 f(y)$ and
\[
Q_{\rm Gr}[f]=\langle F\,\Theta F\rangle
=\int\!\!\int f^\dagger(x)\,
\langle\psi(x)\bar\psi(\Theta y)\rangle\,\gamma_0 f(y)
=\int\!\!\int f^\dagger(x)\,S_E(x-\Theta y)\,\gamma_0\,f(y),
\]
with no reordering sign (fields already ordered as in the pairing).
Here $x-\Theta y=(w_x+w_y,\vec x-\vec y)$, $W=w_x+w_y>0$, $s=+$:
\[
S_E(x-\Theta y)\big|_{\vec p}
=\frac{e^{-\omega W}}{2\omega}
\bigl(m+\omega\gamma_0-i\vec\gamma\cdot\vec p\bigr)
\;\Longrightarrow\;
\mathcal M_{\rm I}(\vec p)
=\frac{(m+\omega\gamma_0-i\vec\gamma\cdot\vec p)\gamma_0}{2\omega}
=\frac{\omega+h_+(\vec p)}{2\omega}.
\]
$\mathcal M_{\rm I}$ is an orthogonal projector, hence
\[
Q_{\rm Gr}[f]
=\int\frac{d^3p}{(2\pi)^3}\,
h_f(\vec p)^\dagger\,\mathcal M_{\rm I}(\vec p)\,h_f(\vec p)\;\ge0,
\qquad
h_f(\vec p)=\int_0^\infty\!dw\,e^{-\omega w}\tilde f(w,\vec p).
\]
\emph{Family II.}\;
$G=\bar\Psi[g]\equiv\int d^4x\,\bar\psi(x)\,g(x)$.
Then $\Theta G=\int g^\dagger(y)\,\gamma_0\,\psi(\Theta y)$ and
\[
Q_{\rm Gr}[g]=\langle G\,\Theta G\rangle
=\int\!\!\int g^\dagger(y)\gamma_0\,
\langle\bar\psi_\alpha(x)\psi_\beta(\Theta y)\rangle
\big|_{\text{matrix}}\,g(x)
=\sigma\int\!\!\int g^\dagger(y)\,\gamma_0\,
S_E(\Theta y-x)\,g(x),
\]
where the reordering rule of \S0 was used once
($\sigma=-1$ here).
Now $\Theta y-x=(-W,\vec y-\vec x)$, $s=-$:
\[
\gamma_0\,S_E(\Theta y-x)\big|_{\vec p}
=\frac{e^{-\omega W}}{2\omega}\,
\gamma_0\bigl(m-\omega\gamma_0-i\vec\gamma\cdot\vec p\bigr)
=\frac{e^{-\omega W}}{2\omega}\,
\bigl(m\gamma_0-\omega-i\gamma_0\vec\gamma\cdot\vec p\bigr)
=e^{-\omega W}\,\frac{h_-(\vec p)-\omega}{2\omega},
\]
and with $\sigma=-1$:
\[
\mathcal M_{\rm II}(\vec p)
=(-1)\cdot\frac{h_-(\vec p)-\omega}{2\omega}
=\frac{\omega-h_-(\vec p)}{2\omega}\;\ge0
\quad(\text{orthogonal projector}).
\]
Hence $Q_{\rm Gr}[g]\ge0$ as well.
[The spatial-momentum orientation differs between the two families
($h_+$ vs $h_-$, equivalently $\vec p\to-\vec p$); both matrices
are projectors, so positivity is orientation-independent.]
\emph{Conclusion (direct direction).}
With the convention set of \S0, both elementary quadratic families
are reflection positive; the extension to the full graded
polynomial algebra is the standard OS fermionic construction and is
cited, not re-proven, here.

% ------------------------------------------------------------
\subsection*{3. Commuting case: explicit negative modes}
The strict negative-mode test below assumes $m>0$; the massless
case requires a separate infrared treatment and is not part of
Q11.1a.
Repeat Family II verbatim with commuting spinor fields:
the \emph{only} change is $\sigma=+1$ in the reordering step, so
\[
\mathcal M_{\rm II}^{\rm com}(\vec p)
=\frac{h_-(\vec p)-\omega}{2\omega},
\qquad
\mathrm{spec}\,\mathcal M_{\rm II}^{\rm com}=\{0,\,-1\}
\quad\text{(each twice)} .
\]
The OS form is not merely indefinite: it is $\le0$ on the whole
Family II and equals $-1$ on the $h_-=-\omega$ eigenmodes.
\emph{Explicit negative-norm test function.}
Take $\vec p_\ast=0$ (envelope concentrated near zero spatial
momentum), a spinor $u_-$ with $\gamma_0u_-=-u_-$ (then
$h_-(0)u_-=-mu_-=-\omega u_-$), and
$g(w,\vec x)=\chi(w)\rho(\vec x)\,u_-$, $\chi$ smooth with support
in $w>0$, $\rho$ a broad envelope.
Then
\[
Q_{\rm com}[g]\;\longrightarrow\;
-\,\bigl|h_g(0)\bigr|^2\;<\;0 .
\]
Family I remains positive in the commuting case (no reordering
occurs there); reflection positivity nevertheless fails because the
OS form must be non-negative on \emph{all} test functionals.
\emph{Mechanism (exact statement).}
The statistics assignment enters the OS test in exactly one place:
the reordering sign $\sigma$ of the
$\langle\bar\psi\,\psi\rangle$ pairing.
Grassmann statistics ($\sigma=-1$) converts the Family-II kernel
into the orthogonal projector $(\omega-h_-)/2\omega$; commuting
statistics ($\sigma=+1$) leaves the negative matrix
$(h_--\omega)/2\omega$.
This replaces the v1 ``symmetrised-kernel'' shortcut: no
charge-conjugation or transpose operation is needed at this level,
because both families use the natural $\psi\bar\psi$ index order
throughout.

\emph{Diagnostic remark (former ``second horn''; not a theorem).}
A separate consistency route would require constructing a positive
commuting Gaussian covariance with spinor indices.
If enforcing ordinary symmetric covariance removes the first-order
Dirac numerator, the resulting object is no longer the same minimal
Dirac kernel.
This is a useful diagnostic, but the converse direction above rests
on the explicit OS negative-mode test, not on this remark.

% ------------------------------------------------------------
\subsection*{4. Caveats and verification list}
\begin{enumerate}[label=(\alph*)]
  \item The trial reflection matrix $\gamma_0$ and the order
    convention $(F,G)=\langle F\,\Theta G\rangle$ must be checked
    against the standard Euclidean fermion OS literature
    (Osterwalder--Schrader Fermi fields; L\"uscher's transfer-matrix
    construction; Montvay--M\"unster \S4) and against the Euclidean
    adjoint/charge-conjugation conventions ultimately adopted in
    the main text.
    The v1 sign mismatch arose precisely from leaving this order
    unfixed; v2 fixes it and derives both family kernels from it.
  \item Spatial-momentum orientation bookkeeping
    ($h_+$ vs $h_-$ between the families) should be re-verified
    once explicit main-text conventions exist; positivity
    conclusions are orientation-independent (projector property),
    the negative-mode conclusion uses only
    $\mathrm{spec}(h_-)=\{\pm\omega\}$.
  \item The extension from the two elementary quadratic families to
    the full graded polynomial algebra is cited from the standard
    OS fermionic construction, not re-derived.
  \item Scope: minimal first-order kernel only; $m>0$; pointlike
    primary-branch excitations; no higher-derivative kernels.
  \item Euclidean gamma conventions are still absent from the main
    text (\S sec:dirac\_eq is Lorentzian); adoption required at
    implementation time.
\end{enumerate}
% ============================================================
